\subsection{Using Verbatim Environment Contents}
\begin{function}{\enverbBody}
  \begin{syntax}
    \cs{enverbBody}
  \end{syntax}
  This macro holds the contents of the last scanned \cs{enverb}-body,
  assignments to it by the package are all local.
\end{function}

\begin{function}{\enverbExecute}
  \begin{syntax}
    \cs{enverbExecute}
  \end{syntax}
  This function executes the current contents of \cs{enverbBody} as though they
  would be part of your input at this place (using |\scantokens|, it is tried as
  good as possible to tackle |\scantokens|'s shortcomings; also it is assumed
  that you didn't change the |eol| key, or else the newlines aren't carried over
  but whatever your |eol| value set is used). The entire contents are
  |\detokenize|d once before |\scantokens| does its magic (shouldn't affect the
  verbatimly read material, but is deemed a sensible safety measure if something
  unexpected happens; this also assumes that anything you placed inside |bol| or
  |eol| is rescanned into the same stuff once |\detokenize|d).\par
  The function isn't fully expandable, the contents are executed at the current
  group level.
\end{function}

\begin{function}{\enverbListing}
  \begin{syntax}
    \cs{enverbListing} \marg{env} \marg{args}
  \end{syntax}
  This function is meant to reuse the current contents of \cs{enverbBody} inside
  another verbatim environment. The used mechanism is tested to work with
  \env{verbatim}, and environments defined by \pkg{fancyvrb}, \pkg{listings} and
  \pkg{minted} (provided you use the default |eol| setting). \meta{env} is the
  environment name in which the contents should be nested, and \meta{args} is
  material to be placed after the |\begin| statement (just as you input them,
  with the outer braces from this argument removed).\par
  Just like |\enverbExecute| this |\detokenize|s the current content of
  \cs{enverbBody} as well as your \meta{args} and uses |\scantokens|.
\end{function}
